\section*{Bab \ref{ch:b2:precambrian-earth} --- Bumi Prakambrium}

\begin{solution}{exo:b2:precambrian-earth:1}
(a) \omterm{def:b2:precambrian-earth:eons}{Hadean} $4567 - 4031 = \qty{536}{Ma}$; \omterm{def:b2:precambrian-earth:eons}{Arkean} $4031 - 2500 = \qty{1531}{Ma}$;
\omterm{def:b2:precambrian-earth:eons}{Proterozoikum} $2500 - 538.8 = \qty{1961}{Ma}$. (b) $(4567 - 538.8)/4567 = 0.882$: sekitar tujuh
perdelapan. (c) $4567 - 4404 = \qty{163}{Ma}$ untuk zirkon itu; $4567 - 4030 = \qty{537}{Ma}$ untuk
batuan tertua.
\end{solution}

\begin{solution}{exo:b2:precambrian-earth:2}
(a) Faktor-faktor $e^{\lambda t}$ dengan $\lambda = \ln2/T$ dan $t = \qty{3}{Ga}$ adalah 1.593
(\ce{^{238}U}), 19.18 (\ce{^{235}U}), 1.160 (\ce{^{232}Th}), dan 5.292 (\ce{^{40}K}): kontribusinya
adalah 3.053, 1.573, 2.424, dan \qty{4.329}{pW}, seluruhnya \qty{11.38}{pW}, dibandingkan dengan
\qty{4.906}{pW} kini. (b) $11.38/4.906 = 2.32$. (c) Torium-232 memimpin kini; kalium-40 memimpin
saat itu.
\end{solution}

\begin{solution}{exo:b2:precambrian-earth:3}
(a) $1616 + 18.0\times(28 - 32.2) = \qty{1540}{\celsius}$. (b) $32.2 + (1500 - 1616)/18.0 =
\qty{25.8}{\percent}$. (c) Analisis batuan utuh bukanlah komposisi cairan: aliran
\omterm{def:b2:precambrian-earth:greenstone}{komatiit} dapat mengandung kristal olivin yang menumpuk dari cairan, yang menambah MgO. Nilai
\qty{38}{\percent}-nya terletak di luar \qty{32}{\percent} cairan \omterm{def:b2:precambrian-earth:greenstone}{komatiit} terkaya, dan
di luar eksperimen yang menjadi dasar garis itu.
\end{solution}

\begin{solution}{exo:b2:precambrian-earth:4}
Batuan tertua, \qty{4.03}{Ga}; stromatolit pertama, \qty{3.43}{Ga}; akhir kubah-kubah Pilbara,
\qty{3.1}{Ga}; \omterm{def:b2:precambrian-earth:goe}{Peristiwa Oksidasi Besar}, sekitar 2.4 sampai \qty{2.1}{Ga}; \omterm{def:b2:precambrian-earth:greenstone}{komatiit} umum terakhir,
\qty{1.9}{Ga}; kenaikan oksigen kedua, sekitar \qty{0.8}{Ga}; akhir Marinoan,
\qty{0.635}{Ga}; fosil Ediakaran pertama, sekitar \qty{0.57}{Ga}.
\end{solution}

\begin{solution}{exo:b2:precambrian-earth:5}
(a) $\tfrac14\times1361\times0.706 = \qty{240.2}{W.m^{-2}}$ diserap dan $\sigma\times288^4 =
\qty{390.1}{W.m^{-2}}$, sehingga $\varepsilon = 240.2/390.1 = 0.616$. (b) $\tfrac14\times1361\times0.4
= \qty{136.1}{W.m^{-2}}$, $T^4 = 136.1/(0.616\times\num{5.670e-8})$, $T = \qty{249.8}{K}$.
(c) Pada \qty{275.5}{K} albedonya adalah $0.6 - 0.306\times12.5/20 = 0.4088$: $R = 340.25\times0.5912 =
\qty{201.2}{W.m^{-2}}$ dan $E = 0.616\,\sigma\times275.5^4 = \qty{201.2}{W.m^{-2}}$. (d) $R' =
340.25\times0.306/20 = \qty{5.21}{W.m^{-2}.K^{-1}}$, $E' = 4\varepsilon\sigma T^3 =
\qty{2.92}{W.m^{-2}.K^{-1}}$: $R' > E'$, sehingga kesetimbangan itu tidak stabil
(\cref{prop:b2:precambrian-earth:ebm}).
\end{solution}

\begin{solution}{exo:b2:precambrian-earth:6}
(a) $1/(1 + 0.4\times2.4/4.567) = 0.826$; $1/(1 + 0.4\times0.65/4.567) = 0.946$. (b) Keadaan hangat
memerlukan sedikitnya 0.932 sinar matahari kini: tidak mungkin \qty{2.4}{Ga} yang lalu, mungkin
\qty{0.65}{Ga} yang lalu. (c) $0.826/0.946 = 0.873$: \qty{13}{\percent} lebih redup, dekat dengan
\qty{10}{\percent} yang dikutip, yang berasal dari model Matahari yang lain.
\end{solution}

\begin{solution}{exo:b2:precambrian-earth:7}
(a) Sturtian $716.9 - 658.9 = \qty{58.0}{Ma}$; Marinoan $644.5 - 635.4 = \qty{9.1}{Ma}$;
masa antarglasial $658.9 - 644.5 = \qty{14.4}{Ma}$. (b) $(58.0 + 9.1)/85 = 0.79$: sekitar empat perlima
periode itu. (c) Awalnya hanya diapit antara 649.9 dan \qty{639.0}{Ma}: menurut batas-batas
itu Marinoan berlangsung antara sekitar 3 dan \qty{15}{Ma}.
\end{solution}

\begin{solution}{exo:b2:precambrian-earth:8}
(a) $\tau = \frac{0.05}{0.95}\times\frac{44.01}{28.97}\times\frac{\num{5.1e18}}{\num{0.26e12}} =
0.05263\times1.519\times\num{1.962e7} = \qty{1.57}{Ma}$. (b) $0.05/\num{278e-6} = 180$ kali.
(c) Dua kali lebih lama, \qty{3.14}{Ma}.
\end{solution}

\begin{solution}{exo:b2:precambrian-earth:9}
(a) Batu-batu itu jatuh ke dalam lumpur lunak dan melengkungkan laminanya ke bawah: batu-batu itu adalah batu jatuhan, yang dilepaskan oleh
es terapung. (b) Es yang mencapai laut \ang{10} dari khatulistiwa adalah jejak yang diharapkan dari sebuah
bola salju, dan sangat sulit diperoleh jika hanya daerah kutub yang membeku. (c) \omterm{def:b2:precambrian-earth:snowball}{Karbonat penudung}:
batuan itu merekam akhir glasiasi yang mendadak, dengan pelapukan yang hebat dan banjirnya
paparan oleh air lelehan.
\end{solution}

\begin{solution}{exo:b2:precambrian-earth:10}
(a) $0.94 > 0.932$, titik kritis hangat: kesetimbangan hangat, $288\times0.94^{1/4} =
\qty{283.6}{K}$, terletak tepat di atas \qty{283}{K}; dan $0.94 < 1.227$, sehingga bola salju juga ada.
(b) $249.8\times0.94^{1/4} = \qty{246.0}{K}$. (c) $\varepsilon = 0.94\times1361\times0.4/(4\sigma
\times263^4) = 511.7/1085.1 = 0.472$. (d) $\sigma\times263^4 = \qty{271.3}{W.m^{-2}}$, dan $(0.616
- 0.472)\times271.3 = \qty{39}{W.m^{-2}}$.
\end{solution}

\begin{solution}{exo:b2:precambrian-earth:11}
(a) $S_{\text{m}}/S_{\text{f}} = (T_1/T_2)^4\,(1 - a_{\text{w}})/(1 - a_{\text{i}})$:
emisivitasnya saling meniadakan. $(263/283)^4\times0.706/0.4 = 0.7459\times1.765 = 1.317$. (b) Kontras
albedo yang lebih kecil antara es dan air terbuka, misalnya es yang lebih gelap dan berdebu, atau lereng
suhu yang lebih lebar. (c) Bola salju memantulkan begitu banyak sinar matahari sehingga tetap dingin pada sinar matahari atau
efek rumah kaca mana pun yang akan membuat Bumi hangat membeku; untuk mencairkannya, gaya pemaksa harus melampaui gaya pemaksa
pembekuan sebesar seluruh lebar simpal itu.
\end{solution}

\begin{solution}{exo:b2:precambrian-earth:12}
Di dekat $T^*$, $C\dot\delta = -(E' - R')\delta$, sehingga $\delta \propto e^{-t/\tau}$ dengan $\tau =
C/(E' - R')$. Pada kesetimbangan hangat $R' = 0$ dan $E' = 4\times0.616\times\num{5.670e-8}\times
288^3 = \qty{3.34}{W.m^{-2}.K^{-1}}$; $C = 1025\times3992\times50 = \qty{2.05e8}{J.m^{-2}.K^{-1}}$,
sehingga $\tau = \qty{6.1e7}{s}$, sekitar \qty{1.9}{yr}. Dalam model ini kesetimbangan hangat mempertahankan $R' =
0$ sampai mencapai \qty{283}{K}, sehingga $\tau$ tetap sekitar dua tahun sampai titik kritis, lalu
keadaan itu lenyap. Dengan albedo yang berubah secara mulus, $R'$ tumbuh selagi kesetimbangan mendekati
lipatan, tempat $E' = R'$: $\tau$ tumbuh tanpa batas, suatu perlambatan yang menandakan titik
kritis.
\end{solution}

\begin{solution}{pb:b2:precambrian-earth:1}
\textbf{1.} $\num{0.0203e-6}\times0.9929\times\num{95.13e-6} = \qty{1.92}{pW.kg^{-1}}$
(\ce{^{238}U}); $\num{0.0203e-6}\times0.00711\times\num{568.47e-6} = \qty{0.0820}{pW.kg^{-1}}$
(\ce{^{235}U}); $\num{0.0795e-6}\times\num{26.28e-6} = \qty{2.09}{pW.kg^{-1}}$ (\ce{^{232}Th});
$\num{240e-6}\times\num{1.20e-4}\times\num{28.47e-6} = \qty{0.820}{pW.kg^{-1}}$ (\ce{^{40}K}).
\textbf{2.} \qty{4.91}{pW.kg^{-1}}.
\textbf{3.} $\ln2/T$: 0.155, 0.985, 0.0495, dan \qty{0.555}{Ga^{-1}}.
\textbf{4.} Faktor 1.52, 14.3, 1.14, dan 4.48: 2.91, 1.17, 2.39, dan \qty{3.67}{pW.kg^{-1}}.
\textbf{5.} \qty{10.1}{pW.kg^{-1}}, 2.07 kali total kini.
\textbf{6.} Kalium-40. Waktu paruhnya, \qty{1.25}{Ga}, pendek: lebih dari dua waktu paruh
telah berlalu sejak itu, dan lebih dari tiga perempat \ce{^{40}K} masa itu telah meluruh.
\textbf{7.} $1/(1 + 0.4\times2.7/4.567) = 0.809$; $0.809\times1361 = \qty{1101}{W.m^{-2}}$.
\textbf{8.} $288\times0.809^{1/4} = \qty{273}{K}$, di bawah \qty{283}{K}: bukan keadaan hangat yang sah;
Bumi akan membeku.
\textbf{9.} $(1101\times0.4/(4\times0.616\,\sigma))^{1/4} = \qty{237}{K}$.
\textbf{10.} $0.616\times0.809 = 0.498$.
\textbf{11.} $0.004/\num{278e-6} = 14$ sampai $0.03/\num{278e-6} = 108$ kali.
\textbf{12.} Termostat pelapukan silikat: pelapukan melambat dalam iklim dingin, sehingga
karbon dioksida vulkanik terakumulasi, dan mempercepat dalam iklim hangat.
\textbf{13.} 0.826 sinar matahari kini, di bawah 0.932: tidak ada keadaan hangat; Bumi dengan udara masa kini akan menjadi
bola salju.
\textbf{14.} $\frac{0.1}{0.9}\times1.519\times\num{1.962e7} = \qty{3.31}{Ma}$.
\textbf{15.} $\frac{0.2}{0.8}\times1.519\times\num{1.962e7} = \qty{7.45}{Ma}$.
\textbf{16.} \omterm{def:b2:precambrian-earth:snowball}{Karbonat penudung}, batu dolomit berlaminasi yang menutupi secara tajam endapan glasial,
yang mengendap dari laut yang dipasok oleh pelapukan hebat pada masa sesudah yang panas dan dibanjiri oleh
air lelehan.
\textbf{17.} Besi(II) hanya tetap terlarut dalam air tanpa oksigen: besi terakumulasi di
samudra anoksik di bawah es dan mengendap sebagai besi(III) di tempat oksigen datang, misalnya bersama
air lelehan.
\textbf{18.} $(1616 - 1450)/(32.2 - 23) = \qty{18.0}{\celsius}$ per persen.
\textbf{19.} $1616 + 18.0\times(30 - 32.2) = \qty{1576}{\celsius}$.
\textbf{20.} \qty{1616}{\celsius}, titik yang terukur.
\textbf{21.} $1616 - 18.0\times20.2 = \qty{1252}{\celsius}$.
\textbf{22.} $1616 - 20\times20.2 = \qty{1212}{\celsius}$: \qty{40}{\celsius} lebih rendah, sepersepuluh
selisih antara kedua lava; kesimpulannya tidak bergantung pada pilihan itu.
\textbf{23.} \omterm{def:b2:precambrian-earth:greenstone}{Komatiit} berasal dari plume, yang lebih panas daripada mantel sekitarnya, dan dari derajat pelelehan
yang besar di kedalaman; suhu sebuah lava bukanlah suhu potensial sumbernya.
Mantel secara keseluruhan lebih hangat sebesar 200 sampai \qty{300}{K}.
\textbf{24.} Lava \omterm{def:b2:precambrian-earth:greenstone}{komatiit} terpanas sekitar \textbf{\qty{364}{\celsius}} lebih panas daripada
basal primer punggungan masa kini.
\end{solution}
